\appendix

\appendix
\section{The curvature term in the Weitzenb\"ock formula}\label{appendix:Weitzen}

In this section, we recall a well-known estimate for the curvature term in the Weitzenb\"ock formula. Let us fix integers $n,N \geq 2$. Let $(M,g)$ be a Riemannian spin manifold of dimension~$n$, and let $f\colon (M,g) \to \big(S^N,g_{S^N}\big)$ be a smooth map to the round unit sphere of dimension $N$. Let~${S \to M}$ denote the spinor bundle of $M$, let $E_0 \to S^N$ denote the spinor bundle of $S^N$ and set $E = f^* E_0$. Let $\mathcal{R}^E$ denote the curvature term appearing in the Weitzenb\"ock formula for the square of the twisted Dirac operator on $S \otimes E$, so that
\[
\big(\mathcal{D}^{S \otimes E }\big)^2 s = \big(\nabla^{\tilde{S} \otimes \tilde{E}}\big)^*\nabla^{\tilde{S} \otimes \tilde{E}} + \frac{1}{4} R_g s + \mathcal{R}^{E} s.
\]
The following estimate for the curvature term $\mathcal{R}^E$ is well known.

\begin{Proposition}[cf.\ Llarull \cite{Llarull}]
\label{curvature.term}
Let $x \in M$ and let $ \mu_1,\hdots,\mu_n \geq 0$ denote the singular values of the differential ${\rm d}f_x\colon (T_x M , g_x) \to \big(T_{f(x)} S^N,g_{S^N}\big)$. Then
\[
 \big|\mathcal{R}^E s\big| \leq \frac{1}{4} \sum_{\substack{1 \leq j,k \leq n\\ j \neq k}} \mu_j \mu_k |s|
\]
for all $s \in S_x \otimes E_x$.
\end{Proposition}

\begin{proof}
Let $m$ denote the rank of the differential ${\rm d}f_x$. Clearly, $m \leq \min \{n,N\}$. We assume that the singular values are arranged so that $\mu_k > 0$ for $1 \leq k \leq m$ and $\mu_k = 0$ for $m+1 \leq k \leq N$. We can find an orthonormal basis $\{e_1,\hdots,e_n\}$ of $(T_x M, g_x)$ and an orthonormal basis $\{\varepsilon_1, \ldots, \varepsilon_N\}$ of $\big(T_{f(x)} S^N, g_{S^N}\big)$ such that ${\rm d}f_x(e_k) = \mu_k \varepsilon_k$ for $1 \leq k \leq m$ and ${\rm d}f_x(e_k) = 0$ for $m+1 \leq k \leq n$.

Let $F^{E_0} \in \Omega^2(S^N, {\rm End}(E_0))$ denote the curvature of $E_0$, and let $F^E \in \Omega^2(M, {\rm End}(E))$ denote the curvature of $E$. For each $s \in S_x \otimes E_x = S_x \otimes (E_0)_{f(x)}$, formula~(8.22) in \cite[Chapter II]{Lawson-Michelsohn} gives
\begin{align*}
\mathcal{R}^E s
&= \frac{1}{2} \sum_{\substack{1 \leq j,k \leq n\\ j \neq k}} \big((e_j \cdot e_k) \otimes F^E_{e_j, e_k}\big) \cdot s = \frac{1}{2} \sum_{\substack{1 \leq j,k \leq n\\ j \neq k}} \big((e_j \cdot e_k) \otimes F^{E_0}_{{\rm d}f_x(e_j), {\rm d}f_x(e_k)}\big) \cdot s \\
&= \frac{1}{2} \sum_{\substack{1 \leq j,k \leq m\\ j \neq k}} \mu_j \mu_k \big((e_j \cdot e_k) \otimes F^{E_0}_{\varepsilon_j,\varepsilon_k}\big) \cdot s.
\end{align*}
Since the curvature operator of $S^N$ acts as the identity on $2$-forms, we obtain by formula~(4.37) in \cite[Chapter II]{Lawson-Michelsohn} (also compare \cite[Lemma 4.3]{Llarull}) that $F^{E_0}_{\varepsilon_j,\varepsilon_k} \eta = \frac{1}{2} \varepsilon_k \cdot \varepsilon_j \cdot \eta$ for all $j \neq k$ and all $\eta \in (E_0)_{f(x)}$.
Putting everything together, it follows that
\[
 \mathcal{R}^E s = \frac{1}{4} \sum_{\substack{1 \leq j,k \leq m\\ j \neq k}} \mu_j \mu_k ((e_j \cdot e_k) \otimes (\varepsilon_k \cdot \varepsilon_j)) \cdot s.
\]
For each pair $j \neq k$, Clifford multiplication by $(e_j \cdot e_k) \otimes (\varepsilon_k \cdot \varepsilon_j)$ on $S_x \otimes E_x = S_x \otimes (E_0)_{f(x)}$ is a self-adjoint involution, hence an isometry.
Therefore, $
 |((e_j \cdot e_k) \otimes (\varepsilon_k \cdot \varepsilon_j)) \cdot s| = |s|$ for all $1 \leq j,k \leq m$.
This finally implies
\begin{equation*}
\big|\mathcal{R}^E s\big| \leq \frac{1}{4} \sum_{\substack{1 \leq j,k \leq m\\ j \neq k}} \mu_j \mu_k |s|.\tag*{\qed}
\end{equation*} \renewcommand{\qed}{}
\end{proof}

\section{A holographic index theorem}

\label{holographic}

We prove a theorem which relates the index on a manifold with boundary with that of the boundary.
We only need it for twisted spinorial Dirac operators but since it may be of independent interest, we show it for the larger class of self-adjoint Dirac-type operators.

Let $M$ be a compact Riemannian manifold with boundary $\partial M$ with outward unit normal vector field $\nu$.
Let $S \to M$ be a Hermitian vector bundle.
Let $\mathcal{D}$ be a differential operator of first order.
Its principal symbol is characterized by $\mathcal{D}(fs)=f\mathcal{D}s+\sigma_{\mathcal{D}}({\rm d}f)s$.
We say that $\mathcal{D}$ is of \emph{Dirac type} if its principal symbol satisfies the Clifford relations
\[
\sigma_{\mathcal{D}}(\xi)\sigma_{\mathcal{D}}(\eta) + \sigma_{\mathcal{D}}(\eta)\sigma_{\mathcal{D}}(\xi) = -2g(\xi,\eta)
\]
for all $\xi,\eta\in T^*_xM$ and $x\in M$.
In particular, if $\mathcal{D}$ is of Dirac type, then $\mathcal{D}$ is elliptic.
All generalized Dirac operators in the sense of Gromov and Lawson are of Dirac type.

We assume that the restriction of $S$ to the boundary $\partial M$ splits into two orthogonal subbundles, $S|_{\partial M} = S^+\oplus S^-$.
A first-order differential operator $\mathcal{D}^\partial\colon C^\infty(\partial M, S) \to C^\infty(\partial M, S)$ is called \emph{adapted to} $\mathcal{D}$ if the principal symbols are related by
\begin{equation}
\sigma_{\mathcal{D}^\partial}(\xi)=-\sigma_{\mathcal{D}}\big(\nu^\flat\big)^{-1}\sigma_{\mathcal{D}}(\xi)
\label{eq:boundarycompatibility}
\end{equation}
for all $\xi\in T^*\partial M$.
Here $\nu^\flat$ is the $1$-form metrically related to $\nu$.\footnote{In \cite{Baer-Ballmann} and many other references one works with the inward unit normal rather than the outward pointing one. Then there is no $-$ sign in \eqref{eq:boundarycompatibility} and the roles of $V_<$ and $V_>$ in the proof of Theorem~\ref{thm:HolIndex} get interchanged.}
If $\mathcal{D}$ is of Dirac type then so is $\mathcal{D}^\partial$.

If $\mathcal{D}^\partial$ interchanges the subbundles, i.e., $\mathcal{D}^\partial\colon C^\infty\big(\partial M, S^\pm\big) \to C^\infty\big(\partial M, S^\mp\big)$, then we call $\mathcal{D}^\partial$ an \emph{odd operator}.
We denote by $C^\infty_\pm(M,S)$ the space of all sections $s$ of $S$ which are smooth up to the boundary and satisfy $s(x)\in S^\pm_x$ for all $x\in\partial M$.

\begin{Theorem}\label{thm:HolIndex}
Let $\mathcal{D}$ be a formally self-adjoint Dirac-type operator and let $\mathcal{D}^\partial$ be a formally self-adjoint odd operator adapted to $\mathcal{D}$.
Assume that $\sigma_{\mathcal{D}}\big(\nu^\flat\big)$ preserves the bundles $S^+$ and $S^-$ and anti-commutes with $\mathcal{D}^\partial$.
Then the operators $\mathcal{D}\colon C^\infty_\pm(M,S)\to C^\infty(M,S)$ and $\mathcal{D}^\partial\colon C^\infty\big(\partial M,S^\pm\big)\to C^\infty\big(\partial M,S^\mp\big)$ are Fredholm and their indices satisfy
\begin{align}
&\mathrm{ind}(\mathcal{D}\colon C^\infty_+(M,S) \to C^\infty(M,S)) \nonumber \\
&\qquad=
-\mathrm{ind}(\mathcal{D}\colon C^\infty_-(M,S)\to C^\infty(M,S)) \label{eq:ind1}\\
&\qquad=
\frac{1}{2} \mathrm{ind}\big(\mathcal{D}^\partial\colon C^\infty\big(\partial M,S^+\big)\to C^\infty(\partial M,S^-)\big) \label{eq:ind2}\\
&\qquad=
-\frac{1}{2} \mathrm{ind}\big(\mathcal{D}^\partial\colon C^\infty(\partial M,S^-)\to C^\infty\big(\partial M,S^+\big)\big) .\label{eq:ind3}
\end{align}
\end{Theorem}

\begin{proof}
The Fredholm property of $\mathcal{D}\colon C^\infty_\pm(M,S)\to C^\infty(M,S)$ follows from Corollary~7.23, Proposition~7.24, and Theorem~7.17 combined with Corollary~8.6 in \cite{Baer-Ballmann}.
Note that the completeness and coercivity at infinity required in \cite{Baer-Ballmann} is automatic in our situation as $M$ is compact.

Since $\sigma_{\mathcal{D}}\big(\nu^\flat\big)$ preserves the bundles $S^+$ and $S^-$, the boundary conditions $s\in S^+$ and $s\in S^-$ are adjoint to each other.
Corollary~8.6 in \cite{Baer-Ballmann} implies
\begin{align*}
\mathrm{ind}(\mathcal{D}\colon C^\infty_\pm(M,S)\to C^\infty(M,S))
={}&
\dim\ker(\mathcal{D}\colon C^\infty_\pm(M,S)\to C^\infty(M,S)) \\
& -
\dim\ker(\mathcal{D}\colon C^\infty_\mp(M,S)\to C^\infty(M,S)).
\end{align*}
In particular, this proves \eqref{eq:ind1}.

Elliptic operators on compact manifolds without boundary are always Fredholm.
Since $\mathcal{D}^\partial$ is elliptic and formally self-adjoint, its $L^2$-closure is self-adjoint.
Since the $L^2$-closures of the restrictions $C^\infty\big(\partial M,S^\pm\big)\to C^\infty\big(\partial M,S^\mp\big)$ are adjoints of each other, we obtain \eqref{eq:ind3}.
It remains to prove \eqref{eq:ind2}.

Let $V_>$ and $V_<$ denote the subspaces of the Sobolev space $H^{\frac{1}{2}}(\partial M,S)$ spanned by the eigenspaces of $\mathcal{D}^\partial$ corresponding to the positive or negative eigenvalues, respectively.
Let $H\subset C^\infty(\partial M,S)$ denote the kernel of $\mathcal{D}^\partial$.
Since $\mathcal{D}^\partial$ interchanges $S^+$ and $S^-$, we may decompose $H=H^+\oplus H^-$, where $H^\pm := H\cap C^\infty\big(\partial M,S^\pm\big)$.
This gives an $L^2$-orthogonal decomposition
\[
H^{\frac{1}{2}}(\partial M,S) = V_> \oplus H^+ \oplus H^- \oplus V_< .
\]
Let $\sigma$ denote the self-adjoint bundle involution on $S$ with the property that $S^\pm$ are the eigen\-spaces to the eigenvalues $\pm 1$.
Since $\mathcal{D}^\partial$ interchanges $S^+$ and $S^-$, it anti-commutes with $\sigma$.
Thus, $\sigma$~maps the eigenspace of $\mathcal{D}^\partial$ for the eigenvalue $\lambda$ isomorphically onto that of $-\lambda$.

Now any $s\in V_>\oplus V_<$ can be expanded into eigensections, $s=\sum_{\lambda\neq0}s_\lambda$.
If furthermore \smash{$s\in H^{\frac{1}{2}}\big(\partial M,S^+\big)$}, then
\[
\sum_{\lambda\neq0}s_\lambda
=
s
=
\sigma s
=
\sum_{\lambda\neq0}\sigma s_\lambda
\]
and therefore $
\sigma s_\lambda = s_{-\lambda}$. Thus, $(V_>\oplus V_<) \cap H^{\frac{1}{2}}\big(\partial M,S^+\big)$ is the graph of the map $\sigma\colon V_> \to V_<$.
We introduce a deformation parameter $t\in [0,1]$ and consider the graph of $t\sigma$.
More precisely, we put $B_t := H^+ \oplus \{s+t\sigma s \colon s \in V_>\}$.
Each $B_t$ is an $\infty$-regular elliptic boundary condition for~$\mathcal{D}$ in the sense of \cite{Baer-Ballmann}.
By deformation invariance of the index, $\mathrm{ind}(\mathcal{D},B_1) = \mathrm{ind}(\mathcal{D},B_0)$. In other words, $\mathrm{ind}(\mathcal{D}\colon C^\infty_+(M,S)\to C^\infty(M,S))$ coincides with the index of $\mathcal{D}$ subject to the boundary condition $H^+\oplus V_>$.

We observe that $H^+\oplus V_>$ is a finite-dimensional modification of the Atiyah--Patodi--Singer boundary condition $V_>$.
Since $\sigma\big(\nu^\flat\big)$ anti-commutes with $\mathcal{D}^\partial$, the adjoint boundary condition of $V_>$ is \smash{$\big(\sigma\big(\nu^\flat\big)V_>\big)^\perp = V_<^\perp=H\oplus V_>$} where $\perp$ denotes the $L^2$-orthogonal complement in $H^{\frac{1}{2}}(\partial M,S)$. Therefore,
\begin{equation} \label{eq:index_a}
\mathrm{ind}(\mathcal{D},V_>) = -\mathrm{ind}(\mathcal{D},H\oplus V_>).
\end{equation}
By \cite[Corollary~8.8]{Baer-Ballmann},
\begin{equation} \label{eq:index_b}
\mathrm{ind}(\mathcal{D},H\oplus V_>) = \mathrm{ind}(\mathcal{D},V_>)+\dim(H).
\end{equation}
Combining \eqref{eq:index_a} and \eqref{eq:index_b} gives
\[
\mathrm{ind}(\mathcal{D},V_>)
=
-\frac{1}{2} \dim(H).
\]
Using \cite[Corollary~8.8]{Baer-Ballmann} again, we obtain
\begin{align*}
\mathrm{ind}(\mathcal{D}\colon C^\infty_+(M,S)\to C^\infty(M,S))
={}&
\mathrm{ind}\big(\mathcal{D},H^+\oplus V_>\big) =
\mathrm{ind}(\mathcal{D},V_>)+\dim\big(H^+\big) \\
={}&
\frac{1}{2} \big(\dim\big(H^+\big)-\dim(H^-)\big) \\
={}&
\frac{1}{2} \big(\dim\ker\big(\mathcal{D}^\partial\colon C^\infty\big(\partial M,S^+\big)\to C^\infty(\partial M,S^-)\big) \\
& -\dim\ker\big(\mathcal{D}^\partial\colon C^\infty(\partial M,S^-)\to C^\infty\big(\partial M,S^+\big)\big)\big) \\
={}&
\frac{1}{2} \mathrm{ind}\big(\mathcal{D}^\partial\colon C^\infty\big(\partial M,S^+\big)\to C^\infty(\partial M,S^-)\big) .
\end{align*}
This concludes the proof.
\end{proof}

The following consequence is known as cobordism invariance of the index:

\begin{Corollary}\label{cor:CobInv}
In addition to the assumptions in Theorem~{\rm\ref{thm:HolIndex}} assume that $S^\pm$ are the eigensubbundles of $S|_{\partial M}$ of the involution ${\rm i}\sigma_{\mathcal{D}}\big(\nu^\flat\big)$ for the eigenvalues $\pm 1$.
Then $S|_{\partial M} = S^+ \oplus S^-$ and the indices occurring in Theorem~{\rm\ref{thm:HolIndex}} vanish.
\end{Corollary}

\begin{proof}
Without loss of generality we can assume that $M$ is connected.
If $\partial M=\varnothing$, then the assertion is obvious.
Therefore, we assume $\partial M\neq \varnothing$.

We claim that the operators $\mathcal{D}\colon C^\infty_\pm(M,S)\to C^\infty(M,S)$ have trivial kernel.
To see this, suppose that $s\in C^\infty_\pm(M,S)\to C^\infty(M,S)$ satisfies $\mathcal{D}s=0$.
Since $\mathcal{D}$ is formally self-adjoint, we obtain
\begin{align*}
0
&=
\int_M \langle \mathcal{D}s,s\rangle - \int_M \langle s, \mathcal{D}s\rangle
=
\int_{\partial M} \big\langle \sigma_{\mathcal{D}}\big(\nu^\flat\big)s,s\big\rangle
=
\mp {\rm i}\int_{\partial M} |s|^2 .
\end{align*}
Hence, $s|_{\partial M}=0$.

We show that this implies $s=0$ on all of $M$.
By adding a small collar neighborhood to $M$ along $\partial M$ we embed $M$ into an open manifold $\tilde M$.
We extend the bundle $S$ and the Dirac-type operator $\mathcal{D}$ to $\tilde M$.
We extend $s$ by zero to $\tilde M$ and obtain a continuous section $\tilde s$.
Let $\phi$ be a~compactly support test section on $\tilde M$.
Then
\begin{align*}
\int_{\tilde M}\langle \tilde s,\mathcal{D}\phi\rangle
&=
\int_M\langle s,\mathcal{D}\phi\rangle
=
\int_M\langle \mathcal{D}s,\phi\rangle + \int_{\partial M} \big\langle s,\sigma_{\mathcal{D}}\big(\nu^\flat\big)\phi\big\rangle
=
0.
\end{align*}
This shows $\mathcal{D}\tilde s=0$ in the weak sense.
By elliptic regularity theory, $\tilde s$ is smooth and $\mathcal{D}\tilde s=0$ holds classically.
Now $\tilde s$ vanishes on a nonempty open subset of $\tilde M$, $\tilde M$ is connected and Dirac-type operators have the unique continuation property, see, e.g., \cite[Theorem~8.2]{Booss-Wojciechowski}.
Thus, $\tilde s=0$ on all of $\tilde M$.
\end{proof}

We generalize Freed's Theorem~B in \cite{Freed} to Dirac-type operators.

\begin{Corollary}\label{cor:Freed}
Let $\mathcal{D}$, $\mathcal{D}^\partial$, $S$, $M$, and $\nu$ be as in Theorem~{\rm\ref{thm:HolIndex}}.
Let $N_1,\hdots,N_k$ denote the connected components of $\partial M$.
Suppose that $\varepsilon_1,\hdots,\varepsilon_k$ is a collection of integers in $\{1,-1\}$. We define a bundle $S^+$ over $\partial M$ so that the fiber of $S^+$ at a point $x \in N_j$ is the eigenspace of~\smash{${\rm i}\sigma_{\mathcal{D}}\big(\nu^\flat\big)$} with eigenvalue $\varepsilon_j$. Similarly, we define a bundle $S^-$ over $\partial M$ so that the fiber of $S^-$ at a point~${x \in \partial M}$ is the eigenspace of \smash{${\rm i}\sigma_{\mathcal{D}}\big(\nu^\flat\big)$} with eigenvalue $-\varepsilon_j$. Then $S|_{\partial M} = S^+ \oplus S^-$. Moreover,
\begin{align*}
\mathrm{ind}(\mathcal{D}\colon C^\infty_+(M,S)\to C^\infty(M,S))
&=
\sum_{\varepsilon_j=1} \mathrm{ind}\big(\mathcal{D}^\partial\colon C^\infty\big(N_j,S^+\big)\to C^\infty(N_j,S^-)\big) .
\end{align*}
\end{Corollary}

\begin{proof}
By Corollary~\ref{cor:CobInv},
\begin{align*}
\begin{aligned}
0
={}&
\sum_{\varepsilon_j=1} \mathrm{ind}\big(\mathcal{D}^\partial\colon C^\infty\big(N_j,S^+\big)\to C^\infty(N_j,S^-)\big) \\
& + \sum_{\varepsilon_j=-1} \mathrm{ind}\big(\mathcal{D}^\partial\colon C^\infty(N_j,S^-)\to C^\infty\big(N_j,S^+\big)\big),
\end{aligned}
\end{align*}
hence
\begin{align*}
0
={}&
\sum_{\varepsilon_j=1} \mathrm{ind}\big(\mathcal{D}^\partial\colon C^\infty\big(N_j,S^+\big)\to C^\infty(N_j,S^-)\big) \\
& - \sum_{\varepsilon_j=-1} \mathrm{ind}\big(\mathcal{D}^\partial\colon C^\infty\big(N_j,S^+\big)\to C^\infty(N_j,S^-)\big) .
\end{align*}
Therefore,
\begin{align*}
\mathrm{ind}\big(\mathcal{D}^\partial\colon C^\infty\big(\partial M,S^+\big)\to C^\infty(\partial M,S^-)\big)
&=
\sum_j \mathrm{ind}\big(\mathcal{D}^\partial\colon C^\infty\big(N_j,S^+\big)\to C^\infty(N_j,S^-)\big) \\
&=
2 \sum_{\varepsilon_j=1} \mathrm{ind}\big(\mathcal{D}^\partial\colon C^\infty\big(N_j,S^+\big)\to C^\infty(N_j,S^-)\big).
\end{align*}
The result now follows from Theorem~\ref{thm:HolIndex}.
\end{proof}
